\begin{proof}
\emph{(i).} By Fact \claimref{synthesis.imported-local-covariance-with-scaling-of-time-ordered-products}(iii) with $\omega = 2$, the coefficient of the derivatives of order $j$ in $K$ is a polynomial with constant coefficients in $\mathcal{R}$, its covariant derivatives and $m$, of dimension $2 - j$, where $\mathcal{R}$ has dimension $2$ and each covariant derivative and $m$ dimension $1$. The construction is covariant under the rotation gauge (Definition \claimref{obj.synthesis.gtg-flat-spacetime-gauge-fields}(e)), under which the slot vectors $a$, $b$, the derivative directions and $\mathcal{R}$ rotate together and the constants do not change; so each coefficient is a multilinear function of its vector slots, and of $\mathcal{R}$ where it occurs, invariant under rotors, hence annihilated by the action of every bivector $J$ (this is the reading made here of ``local and covariant'' in the Fact, which is cited from memory and sourced for scalar fields), $v \mapsto J\times v$ on vectors and $M \mapsto J\times M(\cdot) - M(J\times\cdot)$ on maps of bivectors. The bivectors $\gamma_{0}\gamma_{1}$, $\gamma_{1}\gamma_{2}$, $\gamma_{2}\gamma_{3}$ generate the Lorentz algebra under the commutator product, so invariance under these three suffices; the check computes the dimension of each invariant space as the corank of the three generators acting on the tensor product, exactly.

\emph{Order $2$:} dimension $0$, so the coefficient is an invariant of four vector slots $(a, b; u, v)$, $u, v$ the derivative directions. The space has dimension $4$; the part symmetric in $(u, v)$ has dimension $2$, spanned by $(a\cdot b)(u\cdot v)$ and $(a\cdot u)(b\cdot v) + (a\cdot v)(b\cdot u)$, which give $K_{1}$ and $K_{2}$ up to the order of the derivatives; the antisymmetric part has dimension $2$ and multiplies $[u\cdot D, v\cdot D]$, which on a torsion-free field is the action of the field strength, a term of order $0$ linear in $\mathcal{R}$. Changing the order of two covariant derivatives in $K_{1}$ or $K_{2}$ likewise adds terms of order $0$ linear in $\mathcal{R}$, and of order $1$ none, the torsion being zero. \emph{Order $1$:} dimension $1$; the only quantity of dimension $1$ is $m$, and the coefficient would be $m$ times an invariant of three vector slots, of which there is none. \emph{Order $0$:} dimension $2$: $m^{2}$ times the one invariant $a\cdot b$ of two vector slots, giving $K_{3}$, and terms linear in $\mathcal{R}$; derivatives of $\mathcal{R}$ have dimension at least $3$. The bilinear forms in $(a, b)$ linear in a map $M$ of bivectors and invariant form a space of dimension $6$; on symmetric maps $3$; on the maps with $\partial_{a}\wedge M(a\wedge c) = 0$, where $\mathcal{R}$ lies by (PI), dimension $2$. The forms $a\cdot M(b)$ and $(a\cdot b)\,\partial_{c}\cdot M(c)$ are invariant and are independent (the check evaluates them on a map of Ricci type), so they span, giving $K_{5}$ and $K_{4}$. On the maps that in addition have $\partial_{a}\cdot M(a\wedge c) = 0$ the dimension is $0$. The identification of maps with tensors in the check uses the invariant form $\langle AB\rangle_{0}$ on bivectors, which is nondegenerate, so the space of invariant forms on a subspace of maps has the dimension of the invariant tensors in that subspace times the two vector slots. This gives the list. The invariance used is that of the connected group; a parity-odd term would appear as an invariant in these counts and none survives.

\emph{(ii).} For a scalar $\chi$, $D\chi = g^{\mu}\partial_{\mu}\chi$ with $g^{\mu} = \hb^{T}(e^{\mu})$, and $D\wedge D\chi = (D\wedge g^{\mu})\,\partial_{\mu}\chi + g^{\nu}\wedge g^{\mu}\,\partial_{\nu}\partial_{\mu}\chi = 0$, the first term by the torsion equation of the definition, clause (f), the second by symmetry. The identity $D\cdot(D\wedge A) = \nabla^{2}A - D(D\cdot A) + \mathcal{R}(A)$ is the classical Weitzenb\"ock formula for vector fields; in the store's conventions it reads as follows. Both sides are covariant under the rotation gauge, so fix a point and a gauge with $\Omega = 0$ there. Write $D_{k} = \gamma_{k}\cdot D = D_{\hb(\gamma_{k})}$ for an orthonormal physical frame $\{\gamma_{k}\}$ of constant vectors, so that $D = \gamma^{k}D_{k}$ and $D_{k}\gamma_{l} = \Omega(\hb(\gamma_{k}))\times\gamma_{l} = 0$ at the point. There
\[
D\cdot(D\wedge A) = \sum_{k,l}\gamma^{k}\cdot\bigl(\gamma^{l}\wedge D_{k}D_{l}A\bigr) = \sum_{k,l}\bigl[(\gamma^{k}\cdot\gamma^{l})\,D_{k}D_{l}A - \gamma^{l}\,\bigl(\gamma^{k}\cdot D_{k}D_{l}A\bigr)\bigr] .
\]
The first sum is $\sum_{k}Q(\gamma_{k})D_{k}D_{k}A$, which is $\nabla^{2}A$ at the point, the correction term $((\gamma_{k}\cdot D)\gamma_{k})\cdot D A$ of the rough Laplacian vanishing there. In the second, $D_{k}D_{l} = D_{l}D_{k} + [D_{k}, D_{l}]$; $\sum_{k}\gamma^{k}\cdot D_{l}D_{k}A = D_{l}(D\cdot A)$ at the point, so that part is $-D(D\cdot A)$. The commutator of the derivatives along the background vector fields $\hb(\gamma_{k})$, $\hb(\gamma_{l})$ is the field strength on their plane plus the derivative along their Lie bracket, whose physical image is $D_{k}\gamma_{l} - D_{l}\gamma_{k}$ by the torsion equation and vanishes at the point; so by clause (c) of the definition $[D_{k}, D_{l}]A = R(\hb(\gamma_{k}\wedge\gamma_{l}))\times A = -\mathcal{R}(\gamma_{k}\wedge\gamma_{l})\cdot A$. With $(B\cdot A)\cdot c = -(B\cdot c)\cdot A$ for a bivector $B$, $\gamma^{k}\cdot\bigl(\mathcal{R}(\gamma_{k}\wedge\gamma_{l})\cdot A\bigr) = \bigl(\gamma^{k}\cdot\mathcal{R}(\gamma_{k}\wedge\gamma_{l})\bigr)\cdot A = \mathcal{R}(\gamma_{l})\cdot A$, and the remaining term is $+\sum_{l}\gamma^{l}\,(\mathcal{R}(\gamma_{l})\cdot A) = \mathcal{R}(A)$, the Ricci map being symmetric under (PI). The sign of this term is the point of convention; the check verifies the identity exactly on a torsion-free polynomial field with a generic vector field, in the store's sign and not the opposite one, with the library whose curvature sign is that of the definition. Hence $K(D\chi) = 0$ for $K = c\,(K_{1} - K_{2} + K_{5})$, on every torsion-free field.

Conversely let $K = \sum c_{i}K_{i}$ with $K(D\chi) = 0$ for all $\chi$, all torsion-free fields and all masses. The constants do not depend on the field or the mass, so the conditions from any one field and any values of the mass are necessary. The check evaluates $K_{i}(D\chi)$ at a point of one polynomial torsion-free field, with its curvature jet exact, for two cubic scalars $\chi$ and two values of $m^{2}$: the resulting linear system in $(c_{1}, \dots, c_{5})$ has rank $4$ and its solutions are the multiples of $(1, -1, 0, 0, 1)$. With the first half this proves the statement. On the flat field $\mathcal{R} = 0$ and $D = \nabla$, so $K A = c\,(\nabla^{2}A - \nabla(\nabla\cdot A))$; contracted with $a$ and applied to $b\,\delta$ this is $-c\,a\cdot T(\delta)(b)$ with $T(f)(a) = (a\cdot\nabla)\nabla f - a\nabla^{2}f$ the transverse operator of Theorem \claimref{synthesis.one-loop-vacuum-polarisation-extends-without-a-regulator-and-ward-preserving-extensions-differ-only-by-the-renormalisation-constants}, whose clause (iii) finds the transverse differences of degree at most $2$ to be the multiples of this term.

\emph{(iii).} The hypothesis is the degree of divergence $2$ that clause (iv) of Theorem \claimref{synthesis.vacuum-polarisation-kernel-on-a-gravitational-field-has-the-flat-leading-part-and-no-weyl-term-of-degree-four} argues. Under it, Fact \claimref{synthesis.imported-epstein-glaser-extension-theorem-and-ward-preserving-time-ordered-products}(ii) gives that two extensions differ by derivatives of $\delta$ of order at most $2$, and Fact \claimref{synthesis.imported-local-covariance-with-scaling-of-time-ordered-products}(iii) that for local covariant extensions with almost homogeneous scaling the difference is a local term in the sense of the statement; by (i) it is $\sum c_{i}K_{i}$. Whether transverse extensions exist is not shown (clause (iv) is an argument). If two exist and both are transverse in one slot, their difference $\Delta$ satisfies $\Delta(D\chi) = 0$ for every test scalar $\chi$ (transversality in the slot of $b$, integrated by parts onto $\chi$), on every field, and (ii) applies. The constants are those of the Fact, the same on every torsion-free field, so $c$ is the flat constant.

\emph{(iv) (argued).} The map $\chi \mapsto \bar\Pi(D\chi)$ vanishes off the diagonal when the unextended kernel is transverse there, so it is supported on the diagonal; that it is a polynomial in the curvature with constant coefficients and of dimension $3$ (one more than the kernel's local terms, from the extra derivative) is the supposition. The count as in (i): three derivatives on $\chi$ with one free vector slot, the invariants of four vector slots symmetrised give $D(D\cdot D\chi)$ up to lower order; one derivative on $\chi$ with $m^{2}$, $\mathcal{R}_{0}$ or the Ricci map; no derivative on $\chi$ with a vector of dimension $3$ built from the curvature, which is $D\mathcal{R}_{0}$ or the divergence of the Ricci map, equal to half of it by the contracted second Bianchi identity (classical, not stored); two derivatives with $m$ times an invariant of three vector slots, none. Subtracting $\sum_{i\geq 2}c_{i}K_{i}$ from $\bar\Pi$ changes the first four coefficients independently, by the rank computation of (ii). For the fifth, symmetry of the kernel makes $(\chi, \psi) \mapsto \int\psi\,D\cdot\bar\Pi(D\chi)$ symmetric; the first four terms are symmetric after integration by parts, and $\int\chi\,(D\mathcal{R}_{0})\cdot D\psi - \int\psi\,(D\mathcal{R}_{0})\cdot D\chi$ is not zero for $\psi = 1$ near the support of $\chi$ unless $\int\chi\,D\cdot D\mathcal{R}_{0} = 0$ for all $\chi$. This is an argument, not a proof: the supposition on the form of the divergence is the content of the literature's Ward identity on a gravitational field.
\end{proof}
